\documentclass{article}
\usepackage{geometry}
\usepackage{minted}
\usepackage{xcolor}

\geometry{margin=1in}

\title{Chain of Responsibility Pattern}

\begin{document}
    
    \maketitle
    
    
    \section{Definition}
        \url{https://www.geeksforgeeks.org/system-design/chain-responsibility-design-pattern/}
        \url{https://refactoring.guru/design-patterns/chain-of-responsibility}
    
        
        The \textbf{Chain of Responsibility} is a behavioral design pattern.
        It allows multiple objects to handle a request without the sender
        knowing which object will process it.
        
        A request is passed along a chain of handlers.
        Each handler decides:
        
        \begin{itemize}
            \item If it can handle the request, it processes it.
            \item Otherwise, it forwards the request to the next handler.
        \end{itemize}
        
        In simple words:
        
        \begin{quote}
            Either I handle it, or I pass it to the next one.
        \end{quote}
    
    
    \section{Simple Python Example}
        
        Suppose we have three roles approving vacation requests:
        
        \begin{itemize}
            \item Supervisor: up to 1 day
            \item Manager: up to 5 days
            \item CEO: any number of days
        \end{itemize}
        
        \subsection{Implementation}
            
            \begin{minted}[fontsize=\small, linenos]{python}
class Handler:
    def __init__(self, next_handler=None):
        self.next_handler = next_handler

    def handle(self, days):
        if self.next_handler:
            return self.next_handler.handle(days)
        return "No one can handle this request"


class Supervisor(Handler):
    def handle(self, days):
        if days <= 1:
            return "Supervisor approved"
        return super().handle(days)


class Manager(Handler):
    def handle(self, days):
        if days <= 5:
            return "Manager approved"
        return super().handle(days)


class CEO(Handler):
    def handle(self, days):
        return "CEO approved"
            \end{minted}
        
        \subsection{Building the Chain}
            
            \begin{minted}[fontsize=\small, linenos]{python}
chain = Supervisor(
            Manager(
                CEO()
            )
        )

print(chain.handle(1))   # Supervisor approved
print(chain.handle(3))   # Manager approved
print(chain.handle(10))  # CEO approved
            \end{minted}
    
    
    \section{Core Idea}
        
        Each handler:
        
        \begin{itemize}
            \item Holds a reference to the next handler
            \item Decides whether to process the request
            \item Forwards the request if it cannot handle it
        \end{itemize}
    
    
    \section{Compilation Note}
        
        To compile this document, you must enable shell escape:
        
        \begin{minted}[fontsize=\small]{bash}
pdflatex -shell-escape yourfile.tex
        \end{minted}

\end{document}


